\section{Discussion}
\iffalse % LAB-REVIEW: discussion to-do list hidden (2026-11-03)
{\color{cfigtodo}\footnotesize\noindent[fig todo: major discussion points to include]
\begin{itemize}
\item \textbf{Bounded exhaustiveness}: complete up to the instruction bound where Stage~1 reaches UNSAT (so far the STT and ReCon models); SPT's Stage~1 stops at its cap, so SPT is not bounded-exhaustive; Stage~1.5 is non-exhaustive by design.
% STAGE NAMES (2026-11-11): in the item above, Stage~1 is now Stage~1a and Stage~1.5 is Stage~1b (Fig.~pipeline renamed).
\item \textbf{Trust base}: each defense's specification is translated by hand into a speculation contract, protection set and leakage function; SpecLFB, InvisiSpec and CleanupSpec currently reuse STT's model; note the per-defense exceptions in the shared skeleton.
\item \textbf{Transmission mechanisms} are defense-specific and hand-written (JSON criteria); STT and Recon still use the generic load gate.
\item \textbf{Harness}: probes never change defense logic, but the speculation and timing directives change prediction and timing on purpose; runs that do not realize the modeled speculation are invalid, never hits.
\item \textbf{Window semantics}: a branch counts as unresolved until its squash lands (decided 2026-11-08), so a transmission after resolution but before the squash (the race) is a violation.
\item \textbf{Threat models differ} across implementations of the same defense (STT-Spectre vs.\ Recon-STT); results are per implementation.
\item \textbf{Complementary to random testing}: \toolname{} replicates the AMuLeT bugs that break a defense's stated guarantee (A-UV1, A-UV3--A-UV6); the others (MSHR contention, instruction-cache fetch, cleanup timing, the store\,$\to$\,TLB leak) lie outside the stated guarantee or outside our model.
\item \textbf{Generality}: other simulators or RTL; more speculation primitives (store-to-load forwarding done; e.g.\ value prediction next).
\auditnote{claim}{claim, low}{Store-to-load forwarding is not done as a primitive. The STL model, stl\_bypass mode and hooks exist and passed controls on 60 tests, but no STL window checker or STL campaign exists, and the ReCon STLF bug came from ordinary forwarding in a branch-speculation test.  \textbf{PROPOSED TEXT} (edits the sentence just before this note): more speculation primitives (store-to-load forwarding in progress; e.g.\textbackslash{} value prediction next)  \textit{Source:} project log 2026-10-30 stl-spec (60-test controls); project log 2026-10-28 16:30 ('The window checker is proposed, not built'); URGENT\_TODO.md fig (d) DECIDE}
\end{itemize}}
\fi
- Sentence about 

A defense's execution contract conservatively specifies which instructions are speculative, and therefore protected. The space of \emph{real} leakage, however, is the intersection of that contract with the speculation a given microarchitecture actually performs (gem5's \emph{speculation primitives}). This leaves a design choice: do we test a contract against the broad space it specifies, or against its projection onto realizable speculation? Testing the realizable intersection is more charitable to implementers --- it surfaces only bugs that can actually manifest --- but is arguably incomplete, since a correct implementation should uphold the \emph{entire} promised contract. The same tension recurs in whether to require misprediction or to also exercise correctly-predicted but still-speculative windows
\urgentbr
\urgentnote{\textbf{VOCAB}: Here ``still-speculative windows'' means stretches where a correctly predicted branch is still unresolved, not the checker's violation window. Reword rather than rename, e.g.\ ``to also exercise correctly predicted branches while they are still unresolved'', so that ``window'' keeps one meaning.}%
\notebr
, and in whether a conservative window (protecting every instruction after an unresolved branch, whether or not it is data-dependent) must be honored even where no real dataflow path to a transmitter exists. \toolname{} can target either interpretation --- for several defenses we maintain both the raw and the intersected execution contract
\notebr
\auditnote{fact}{fact, low}{Unsupported: every model in the evaluation uses the branch-only execution contract, and no defense has both a raw and an intersected contract in use. Broader STT contracts exist only as a draft awaiting approval (store shadows, for ReCon's STT) or unused (Futuristic). \newline \textbf{PROPOSED TEXT} (replaces the clause just before this note): \textbackslash{}toolname\{\} can target either interpretation (for STT we have also written broader execution contracts, which we leave to future work) \newline \textit{Source:} simspect URGENT\_TODO.md DECIDE ``approve the draft STAGE1\_alloy/models/proposed/STT\_6\_shadows.als''; wiki/design/FUTURISTIC\_MODEL.md.}
\notebr
--- and we default to the realizable intersection, since it is what yields demonstrable, root-caused bugs on the running prototype.

\anna{2. Alloy abstract LILT of the STT load load, 3. the enumerated, annotated veersion 4. a really minimalist version of my pipelineout of running on gem5 with the relevant windows circled and completion check occuring -> bug}

\urgentbr\urgentnote{\textbf{PROPOSED TEXT (checked guarantees and scope).} \toolname{} checks each defense against the guarantee its authors state: each footprint criterion lists only the structures and state that the defense's paper claims to protect, together with the sentence that claims it, and where we test less than a paper claims we say so (SpecLFB's protection after faulting instructions, below). For InvisiSpec, that is the data-cache hierarchy, meaning the occupancy, replacement and coherence state of every level, plus the data TLB (\cite[\S III-B]{yan2018invisispec}). Two of the vulnerabilities AMuLeT reports against InvisiSpec fall outside this statement. Speculative instruction fetches into the L1 instruction cache (A-KV1 in~\cite{amulet}) are excluded by the paper's footnote~2: ``For simplicity, we do not describe the protection of the instruction cache or instruction TLB.'' MSHR contention (A-UV2 in~\cite{amulet}) is excluded explicitly in the security analysis: ``contention on other resources, such as MSHRs or execution units, could slow down the receiver, but such side channels are considered out of scope in this paper'' (\cite[\S VII]{yan2018invisispec}). We therefore treat these as limits of the defense's stated scope rather than implementation bugs, and our checker does not flag them, whereas the speculative L1D eviction (A-UV1), which violates the stated data-cache guarantee, is flagged. The same rule applies to the other footprint defenses. CleanupSpec states its guarantee for the data-cache hierarchy (L1-D and L2/LLC); it acknowledges that cleanup time varies and leaves a constant-time cleanup to future work (\cite{saileshwar2019cleanupspec}, \S 4(b)), so a cleanup-latency difference such as unXpec (A-KV2) is not checked. CleanupSpec's failures to undo a speculative install (A-UV3--A-UV5) do violate its stated guarantee, and our checker flags them. SpecLFB excludes TLB and port-contention channels and loads that hit in the cache~\cite{cheng2024speclfb}. \textit{Sources:} InvisiSpec fn.~2 p3 and \S VII p10 of the original MICRO version; CleanupSpec \S 4(b) p8 and \S 2.1 p3; SpecLFB \S 4 p635 and \S 9 p643; criterion files specs/leakage/\{invisispec,cleanupspec,speclfb\}.json.}

\urgentbr\urgentnote{\textbf{PROPOSED TEXT (STT and store addresses).} AMuLeT also reports that a speculative store with a tainted address installs a D-TLB entry under STT (A-KV3 in~\cite{amulet}), a leak first identified by DOLMA~\cite{loughlin2021dolma}. STT ties its protection of translation state to delaying loads (``for loads, delaying execution implies delaying the TLB lookup'', \cite[footnote~2]{yu2019stt}) and ``does not require tracking taint in any part of the memory system (TLB, caches, or memory)''. Our STT model's leakage function contains load addresses and branch predicates, the instructions STT delays, and not store addresses, so \toolname{} does not generate tests aimed at A-KV3. STT also assumes that stores act only at retirement: ``We assume that stores trigger a cache coherence invalidation only on retirement, or else are defined as transmitters'' (\cite{yu2019stt}, p5), and ``In isolation, a store creates neither explicit or implicit channels because we assume stores are performed at retirement'' (p10). Our model makes the same assumption, so A-KV3, a store that translates its address before retirement, lies outside what we check. \textit{Note:} the first quote's alternative (``or else are defined as transmitters'') lets a reader argue that a store acting early is an STT transmitter; the project's decision (simspect URGENT\_TODO.md, ``store addresses under STT'') keeps A-KV3 out of scope, so word this as our model's assumption, not as STT's exclusion.}

\urgentbr\urgentnote{\textbf{PROPOSED TEXT (structure sizes).} AMuLeT amplifies contention-dependent leaks by testing with smaller microarchitectural structures; in its InvisiSpec campaign it reduces the L1D from 8 ways to 2 and the MSHRs from 256 to 2 (\cite[\S 4.5.1, Table~6]{amulet}). We apply the same idea to every footprint defense and run with a small L1D (1--2\,kB, 2--4 ways). Small structures let a short litmus test reach the full-set and eviction conditions under which these bugs appear, and they do not change the defense logic. All timing and throughput comparisons with AMuLeT instead use exactly AMuLeT's structure sizes and configuration.}

\urgentbr\urgentnote{\textbf{PROPOSED TEXT (cache sizes for the taint defenses).} We do not shrink caches for the taint-tracking defenses (STT, ReCon, SPT): their violation test asks whether a transmitter obtained its data inside the violation window, which does not depend on cache pressure, so smaller caches would amplify nothing. Any new footprint bug found on the small configuration is re-run once at AMuLeT's exact configuration before we report it.}
\urgentbr\urgentnote{\textbf{PROPOSED TEXT (why the pending-squash bug is not reached).} A fix in the upstream STT repository (December 2025) addresses a bug that \toolname{} does not reach: STT delays a tainted branch's squash to hide its direction, but when an older tainted branch and a younger untainted branch mispredict in the same cycle, the untainted branch's squash is delayed too, so its redirect timing reveals the tainted branch's direction~\cite{yu2019stt}. Two things keep it out of reach. First, the observable is the squash signal of an \emph{untainted} branch, which is not a transmitter in our STT specification (load addresses and the predicates of tainted branches); our security predicate calls such a gadget secure, and capturing it needs a squash-signal channel in the model, which we have drafted but not evaluated. Second, even with that channel, a test of our size cannot make the signal depend on the secret: tainting the older branch requires a speculative load, which requires a still older unresolved branch, and gem5 holds both pending squashes until that oldest branch resolves, so the younger branch's redirect happens when the window closes whatever the secret is. Comparing a tainted with an untainted older branch shows the younger branch being held (17 of 20 generated gadgets), but varying only the secret changes nothing observable (0 of 12). Reaching it likely needs the two branches to resolve in the same cycle while the window is opened by something other than an older branch, which we leave to future work.}
\urgentbr\urgentnote{\textbf{PROPOSED TEXT (what we leave out of scope, and why this is closer to each paper).} \toolname{} judges each prototype against the guarantee its own paper states. A report then means that the implementation breaks a promise its authors made, which is a bug they can fix in the code. Checking every defense against a generic leakage contract, as AMuLeT does~\cite{amulet}, instead mixes such bugs with channels the authors chose not to defend, which no change to the implementation would remove. Of the reports we leave out, three kinds follow from the papers or from our specifications. First, channels the defense's paper excludes: InvisiSpec's MSHR contention (A-UV2) and instruction cache (A-KV1), above. Second, the timing of a defense's own actions: CleanupSpec acknowledges that cleanup time varies and leaves a constant-time cleanup to future work (\cite{saileshwar2019cleanupspec}, \S 4(b)); unXpec~\cite{li2022unxpec} shows the variation is exploitable, and AMuLeT observes it (A-KV2) through the instruction cache (\cite{amulet}, p12), which CleanupSpec's threat model excludes (\cite{saileshwar2019cleanupspec}, \S 2.1). Third, transmitters our specifications do not include: speculative stores (STT's TLB leak, A-KV3, above), and the squash signal of an untainted branch (the pending-squash bug, above). \textit{Pending decision} (simspect URGENT\_TODO.md, ``store addresses under STT''): this list used to include the tainted-address store channel in ReCon's implementation (a tainted store address resolves and squashes an aliasing younger load). STT claims to close exactly that channel (its memory-dependence example on p5, and p8: delay squashing ``until the load's address and the addresses of all older stores become untainted''), so it may belong under what we do not yet cover instead; it is left out of both lists until decided. \textit{Check before use:} that ``generic leakage contract'' is a fair description of how AMuLeT chooses its contract per defense.}

\urgentbr\urgentnote{\textbf{PROPOSED TEXT (what we do not yet cover).} Some reports fall within a paper's claim but outside what \toolname{} currently generates; we list them as coverage limits rather than scope. SpecLFB's paper protects loads that follow a possibly faulting instruction, and source reading shows its implementation does not; our speculation contract models only branch and store-address speculation, so \toolname{} cannot yet test this. Source reading also finds candidate bugs that need tests our generator cannot yet produce: interactions between cores (in InvisiSpec, a speculative request that downgrades another core's copy of a line), instructions outside our vocabulary (SPT's untaint rules for flags, conditional moves and 8- and 16-bit writes, and ReCon's floating-point loads, which are never marked speculative), and shapes beyond our bound (a squash in the middle of the window while an older shadow is still open). Each needs an extension of the model or the vocabulary, not of the method.}

\urgentbr\urgentnote{\textbf{PROPOSED TEXT (sketch: why forcing predictions and stretching windows is sound).} The harness changes two things in the simulated core (Speculation Primitive Instrumentation): where the predictors send execution, by forcing each annotated branch to mispredict or to predict correctly and by forcing a load to be predicted not to alias an older store, and when an unresolved branch may resolve or the final instruction may commit, by stalling them for a set number of cycles. (a) This does not change the defense under test. The hooks act only on predictor outputs and on when resolution or commit proceeds; they add prediction overrides, stalls and trace prints, and never modify the defense's own state or code paths (taint, protection flags, speculative buffers, cleanup records). The defense therefore sees an ordinary misprediction and a branch whose condition arrives late, inputs it must already handle. (b) It does not create leakage that a regular predictor could not. In the Spectre threat model the attacker trains the predictors~\cite{spectre}, so any prediction we force is one an attacker can induce, and a late resolution corresponds to a branch condition that misses in the cache, which an attacker can arrange by evicting it. The defenses aim to be secure for any prediction and any speculation window [cite each threat model], so a leak inside a forced, stretched window is one the defense set out to prevent. Two limits remain. A stretched window can make a leak easier to hit than natural timing would, so we confirm new bugs without the hooks where we can (the SpecLFB \texttt{jmp} bug reproduces on an unmodified build with a natural misprediction). Changing a pipeline parameter, such as the delay between commit and issue, changes the microarchitecture itself, so bugs that need it (the visibility-point/squash race) are reported together with that setting.}

\urgentbr\urgentnote{\textbf{PROPOSED TEXT (targets without white-box access; from the old draft, oldtechnicaldetails.tex:42 and :101).} \toolname{} uses gem5's white-box access to force predictions and hold resolution (Gem5 Execution and Instrumentation). A target without that access could realize the same annotations the way an attacker does (argument (b) of the note above): train the predictors before the test so that each annotated branch mispredicts or predicts as its annotation says, and widen the window by evicting the data a branch condition depends on so that the branch resolves late. Detection carries over to RTL simulation, which still exposes every signal; on silicon a hit would have to be measured, for example with a cache-timing probe. \textit{Notes:} none of this is implemented. It is a generality point for the next-steps paragraph or the Conclusion, not a result, and the RTL and silicon sentence is a suggestion, not something we have checked.}
